
        You are a research assistant AI tasked with generating a scientific paper based on provided literature. Follow these steps:
    1. Analyze the given References.
    2. Identify gaps in existing research to establish the motivation for a new study.
    3. Propose a main idea for a new research work.
    4. Write the paper's main content in LaTeX format, including all the following parts, please must have tables:
    - Title
    - Abstract
    - Introduction
    - Related Work
    - Methods/Experiment
    - Results with tables
    - Discussion
    - Conclusion
    - Contributions
    Ensure all content is original, academically rigorous, and follows standard scientific writing conventions.Please make all decision by yourself,do not ask for any instructions

        Here are the references to be used for generating the paper.
        You MUST base the paper on these references and integrate them naturally.

        References:
        --------------------
        @misc{santini2026itsconfidenceunsupervisedapproach,
      title={It's All About the Confidence: An Unsupervised Approach for Multilingual Historical Entity Linking using Large Language Models}, 
      author={Cristian Santini and Marieke Van Erp and Mehwish Alam},
      year={2026},
      eprint={2601.08500},
      archivePrefix={arXiv},
      primaryClass={cs.CL},
      url={https://arxiv.org/abs/2601.08500},
      abstract = {Despite the recent advancements in NLP with the advent of Large Language Models (LLMs), Entity Linking (EL) for historical texts remains challenging due to linguistic variation, noisy inputs, and evolving semantic conventions. Existing solutions either require substantial training data or rely on domain-specific rules that limit scalability. In this paper, we present MHEL-LLaMo (Multilingual Historical Entity Linking with Large Language MOdels), an unsupervised ensemble approach combining a Small Language Model (SLM) and an LLM. MHEL-LLaMo leverages a multilingual bi-encoder (BELA) for candidate retrieval and an instruction-tuned LLM for NIL prediction and candidate selection via prompt chaining. Our system uses SLM's confidence scores to discriminate between easy and hard samples, applying an LLM only for hard cases. This strategy reduces computational costs while preventing hallucinations on straightforward cases. We evaluate MHEL-LLaMo on four established benchmarks in six European languages (English, Finnish, French, German, Italian and Swedish) from the 19th and 20th centuries. Results demonstrate that MHEL-LLaMo outperforms state-of-the-art models without requiring fine-tuning, offering a scalable solution for low-resource historical EL. The implementation of MHEL-LLaMo is available on Github.}
}@misc{finkelstein2026translategemmatechnicalreport,
      title={TranslateGemma Technical Report}, 
      author={Mara Finkelstein and Isaac Caswell and Tobias Domhan and Jan-Thorsten Peter and Juraj Juraska and Parker Riley and Daniel Deutsch and Cole Dilanni and Colin Cherry and Eleftheria Briakou and Elizabeth Nielsen and Jiaming Luo and Kat Black and Ryan Mullins and Sweta Agrawal and Wenda Xu and Erin Kats and Stephane Jaskiewicz and Markus Freitag and David Vilar},
      year={2026},
      eprint={2601.09012},
      archivePrefix={arXiv},
      primaryClass={cs.CL},
      url={https://arxiv.org/abs/2601.09012},
      abstract = {We present TranslateGemma, a suite of open machine translation models based on the Gemma 3 foundation models. To enhance the inherent multilingual capabilities of Gemma 3 for the translation task, we employ a two-stage fine-tuning process. First, supervised fine-tuning is performed using a rich mixture of high-quality large-scale synthetic parallel data generated via state-of-the-art models and human-translated parallel data. This is followed by a reinforcement learning phase, where we optimize translation quality using an ensemble of reward models, including MetricX-QE and AutoMQM, targeting translation quality. We demonstrate the effectiveness of TranslateGemma with human evaluation on the WMT25 test set across 10 language pairs and with automatic evaluation on the WMT24++ benchmark across 55 language pairs. Automatic metrics show consistent and substantial gains over the baseline Gemma 3 models across all sizes. Notably, smaller TranslateGemma models often achieve performance comparable to larger baseline models, offering improved efficiency. We also show that TranslateGemma models retain strong multimodal capabilities, with enhanced performance on the Vistra image translation benchmark. The release of the open TranslateGemma models aims to provide the research community with powerful and adaptable tools for machine translation.}
}@misc{mim2026unimodallanguageagentprovide,
      title={Can a Unimodal Language Agent Provide Preferences to Tune a Multimodal Vision-Language Model?}, 
      author={Sazia Tabasum Mim and Jack Morris and Manish Dhakal and Yanming Xiu and Maria Gorlatova and Yi Ding},
      year={2026},
      eprint={2601.06424},
      archivePrefix={arXiv},
      primaryClass={cs.CL},
      url={https://arxiv.org/abs/2601.06424},
      abstract = {To explore a more scalable path for adding multimodal capabilities to existing LLMs, this paper addresses a fundamental question: Can a unimodal LLM, relying solely on text, reason about its own informational needs and provide effective feedback to optimize a multimodal model? To answer this, we propose a method that enables a language agent to give feedback to a vision-language model (VLM) to adapt text generation to the agent's preferences. Our results from different experiments affirm this hypothesis, showing that LLM preference feedback significantly enhances VLM descriptions. Using our proposed method, we find that the VLM can generate multimodal scene descriptions to help the LLM better understand multimodal context, leading to improvements of maximum 13\% in absolute accuracy compared to the baseline multimodal approach. Furthermore, a human study validated our AI-driven feedback, showing a 64.6\% preference alignment rate between the LLM's choices and human judgments. Extensive experiments provide insights on how and why the method works and its limitations.}
}@misc{vejendla2026rewritenetsendtoendtrainablestringrewriting,
      title={RewriteNets: End-to-End Trainable String-Rewriting for Generative Sequence Modeling}, 
      author={Harshil Vejendla},
      year={2026},
      eprint={2601.07868},
      archivePrefix={arXiv},
      primaryClass={cs.LG},
      url={https://arxiv.org/abs/2601.07868},
      abstract = {Dominant sequence models like the Transformer represent structure implicitly through dense attention weights, incurring quadratic complexity. We propose RewriteNets, a novel neural architecture built on an alternative paradigm: explicit, parallel string rewriting. Each layer in a RewriteNet contains a set of learnable rules. For each position in an input sequence, the layer performs four operations: (1) fuzzy matching of rule patterns, (2) conflict resolution via a differentiable assignment operator to select non-overlapping rewrites, (3) application of the chosen rules to replace input segments with output segments of potentially different lengths, and (4) propagation of untouched tokens. While the discrete assignment of rules is non-differentiable, we employ a straight-through Gumbel-Sinkhorn estimator, enabling stable end-to-end training. We evaluate RewriteNets on algorithmic, compositional, and string manipulation tasks, comparing them against strong LSTM and Transformer baselines. Results show that RewriteNets excel at tasks requiring systematic generalization (achieving 98.7\% accuracy on the SCAN benchmark's length split) and are computationally more efficient than Transformers. We also provide an analysis of learned rules and an extensive ablation study, demonstrating that this architecture presents a promising direction for sequence modeling with explicit structural inductive biases.}
}@misc{meng2026languagemodelsreasonlanguages,
      title={Do Language Models Reason Across Languages?}, 
      author={Yan Meng and Wafaa Mohammed and Christof Monz},
      year={2026},
      eprint={2601.06644},
      archivePrefix={arXiv},
      primaryClass={cs.CL},
      url={https://arxiv.org/abs/2601.06644},
      abstract = {The real-world information sources are inherently multilingual, which naturally raises a question about whether language models can synthesize information across languages. In this paper, we introduce a simple two-hop question answering setting, where answering a question requires making inferences over two multilingual documents. We find that language models are more sensitive to language variation in answer-span documents than in those providing bridging information, despite the equal importance of both documents for answering a question. Under a step-by-step sub-question evaluation, we further show that in up to 33\% of multilingual cases, models fail to infer the bridging information in the first step yet still answer the overall question correctly. This indicates that reasoning in language models, especially in multilingual settings, does not follow a faithful step-by-step decomposition. Subsequently, we show that the absence of reasoning decomposition leads to around 18\% composition failure, where both sub-questions are answered correctly but fail for the final two-hop questions. To mitigate this, we propose a simple three-stage SUBQ prompting method to guide the multi-step reasoning with sub-questions, which boosts accuracy from 10.1\% to 66.5\%.}
}@misc{liang2026hiddenstatesearlysignals,
      title={Hidden States as Early Signals: Step-level Trace Evaluation and Pruning for Efficient Test-Time Scaling}, 
      author={Zhixiang Liang and Beichen Huang and Zheng Wang and Minjia Zhang},
      year={2026},
      eprint={2601.09093},
      archivePrefix={arXiv},
      primaryClass={cs.LG},
      url={https://arxiv.org/abs/2601.09093},
      abstract = {Large Language Models (LLMs) can enhance reasoning capabilities through test-time scaling by generating multiple traces. However, the combination of lengthy reasoning traces with multiple sampling introduces substantial computation and high end-to-end latency. Prior work on accelerating this process has relied on similarity-based or confidence-based pruning, but these signals do not reliably indicate trace quality. To address these limitations, we propose STEP: Step-level Trace Evaluation and Pruning, a novel pruning framework that evaluates reasoning steps using hidden states and dynamically prunes unpromising traces during generation. We train a lightweight step scorer to estimate trace quality, and design a GPU memory-aware pruning strategy that triggers pruning as the GPU memory is saturated by KV cache to reduce end-to-end latency. Experiments across challenging reasoning benchmarks demonstrate that STEP reduces end-to-end inference latency by 45\%-70\% on average compared to self-consistency while also improving reasoning accuracy. Our code is released at: https://github.com/Supercomputing-System-AI-Lab/STEP}
}@misc{ma2026slamllmmodularopensourcemultimodal,
      title={SLAM-LLM: A Modular, Open-Source Multimodal Large Language Model Framework and Best Practice for Speech, Language, Audio and Music Processing}, 
      author={Ziyang Ma and Guanrou Yang and Wenxi Chen and Zhifu Gao and Yexing Du and Xiquan Li and Zhisheng Zheng and Haina Zhu and Jianheng Zhuo and Zheshu Song and Ruiyang Xu and Tiranrui Wang and Yifan Yang and Yanqiao Zhu and Zhikang Niu and Liumeng Xue and Yinghao Ma and Ruibin Yuan and Shiliang Zhang and Kai Yu and Eng Siong Chng and Xie Chen},
      year={2026},
      eprint={2601.09385},
      archivePrefix={arXiv},
      primaryClass={cs.SD},
      doi={https://doi.org/10.1109/JSTSP.2026.3653157},
      url={https://arxiv.org/abs/2601.09385},
      abstract = {The recent surge in open-source Multimodal Large Language Models (MLLM) frameworks, such as LLaVA, provides a convenient kickoff for artificial intelligence developers and researchers. However, most of the MLLM frameworks take vision as the main input modality, and provide limited in-depth support for the modality of speech, audio, and music. This situation hinders the development of audio-language models, and forces researchers to spend a lot of effort on code writing and hyperparameter tuning. We present SLAM-LLM, an open-source deep learning framework designed to train customized MLLMs, focused on speech, language, audio, and music processing. SLAM-LLM provides a modular configuration of different encoders, projectors, LLMs, and parameter-efficient fine-tuning plugins. SLAM-LLM also includes detailed training and inference recipes for mainstream tasks, along with high-performance checkpoints like LLM-based Automatic Speech Recognition (ASR), Automated Audio Captioning (AAC), and Music Captioning (MC). Some of these recipes have already reached or are nearing state-of-the-art performance, and some relevant techniques have also been accepted by academic papers. We hope SLAM-LLM will accelerate iteration, development, data engineering, and model training for researchers. We are committed to continually pushing forward audio-based MLLMs through this open-source framework, and call on the community to contribute to the LLM-based speech, audio and music processing.}
}@misc{torunoğluselamet2026parallelcrosslingualbenchmarkmultimodal,
      title={A Parallel Cross-Lingual Benchmark for Multimodal Idiomaticity Understanding}, 
      author={Dilara Torunoğlu-Selamet and Dogukan Arslan and Rodrigo Wilkens and Wei He and Doruk Eryiğit and Thomas Pickard and Adriana S. Pagano and Aline Villavicencio and Gülşen Eryiğit and Ágnes Abuczki and Aida Cardoso and Alesia Lazarenka and Dina Almassova and Amalia Mendes and Anna Kanellopoulou and Antoni Brosa-Rodríguez and Baiba Saulite and Beata Wojtowicz and Bolette Pedersen and Carlos Manuel Hidalgo-Ternero and Chaya Liebeskind and Danka Jokić and Diego Alves and Eleni Triantafyllidi and Erik Velldal and Fred Philippy and Giedre Valunaite Oleskeviciene and Ieva Rizgeliene and Inguna Skadina and Irina Lobzhanidze and Isabell Stinessen Haugen and Jauza Akbar Krito and Jelena M. Marković and Johanna Monti and Josue Alejandro Sauca and Kaja Dobrovoljc and Kingsley O. Ugwuanyi and Laura Rituma and Lilja Øvrelid and Maha Tufail Agro and Manzura Abjalova and Maria Chatzigrigoriou and María del Mar Sánchez Ramos and Marija Pendevska and Masoumeh Seyyedrezaei and Mehrnoush Shamsfard and Momina Ahsan and Muhammad Ahsan Riaz Khan and Nathalie Carmen Hau Norman and Nilay Erdem Ayyıldız and Nina Hosseini-Kivanani and Noémi Ligeti-Nagy and Numaan Naeem and Olha Kanishcheva and Olha Yatsyshyna and Daniil Orel and Petra Giommarelli and Petya Osenova and Radovan Garabik and Regina E. Semou and Rozane Rebechi and Salsabila Zahirah Pranida and Samia Touileb and Sanni Nimb and Sarfraz Ahmad and Sarvinoz Nematkhonova and Shahar Golan and Shaoxiong Ji and Sopuruchi Christian Aboh and Srdjan Sucur and Stella Markantonatou and Sussi Olsen and Vahide Tajalli and Veronika Lipp and Voula Giouli and Yelda Yeşildal Eraydın and Zahra Saaberi and Zhuohan Xie},
      year={2026},
      eprint={2601.08645},
      archivePrefix={arXiv},
      primaryClass={cs.CL},
      url={https://arxiv.org/abs/2601.08645},
      abstract = {Potentially idiomatic expressions (PIEs) construe meanings inherently tied to the everyday experience of a given language community. As such, they constitute an interesting challenge for assessing the linguistic (and to some extent cultural) capabilities of NLP systems. In this paper, we present XMPIE, a parallel multilingual and multimodal dataset of potentially idiomatic expressions. The dataset, containing 34 languages and over ten thousand items, allows comparative analyses of idiomatic patterns among language-specific realisations and preferences in order to gather insights about shared cultural aspects. This parallel dataset allows to evaluate model performance for a given PIE in different languages and whether idiomatic understanding in one language can be transferred to another. Moreover, the dataset supports the study of PIEs across textual and visual modalities, to measure to what extent PIE understanding in one modality transfers or implies in understanding in another modality (text vs. image). The data was created by language experts, with both textual and visual components crafted under multilingual guidelines, and each PIE is accompanied by five images representing a spectrum from idiomatic to literal meanings, including semantically related and random distractors. The result is a high-quality benchmark for evaluating multilingual and multimodal idiomatic language understanding.}
}@misc{nuraini2026mechanismstransferabledataefficientlowresource,
      title={Mechanisms are Transferable: Data-Efficient Low-Resource Adaptation via Circuit-Targeted Supervised Fine-Tuning}, 
      author={Khumaisa Nur'aini and Ayu Purwarianti and Alham Fikri Aji and Derry Wijaya},
      year={2026},
      eprint={2601.08146},
      archivePrefix={arXiv},
      primaryClass={cs.CL},
      url={https://arxiv.org/abs/2601.08146},
      abstract = {Adapting LLMs to low-resource languages is difficult: labeled data is scarce, full-model fine-tuning is unstable, and continued cross-lingual tuning can cause catastrophic forgetting. We propose Circuit-Targeted Supervised Fine-Tuning (CT-SFT): a counterfactual-free adaptation of CD-T (Contextual Decomposition Transformer) that uses a label-balanced mean baseline and task-directional relevance scoring to identify a sparse set of task-relevant attention heads in a proxy-language checkpoint, then transfer learns to a target language by updating only those heads (plus LayerNorm) via head-level gradient masking. Across NusaX-Senti and XNLI, CT-SFT improves cross-lingual accuracy over continued full fine-tuning while updating only a small subset of model parameters. We find an editing-preserving trade-off: harder transfers favor editing circuit heads, while easier transfers often favor near-zero (i.e., low-relevance heads) updates, preserving the source mechanism. CT-SFT also substantially reduces catastrophic forgetting, preserving proxy/source-language competence during transfer.}
}@misc{tian2026stagebenchmarkknowledgegraph,
      title={STAGE: A Benchmark for Knowledge Graph Construction, Question Answering, and In-Script Role-Playing over Movie Screenplays}, 
      author={Qiuyu Tian and Yiding Li and Fengyi Chen and Zequn Liu and Youyong Kong and Fan Guo and Yuyao Li and Jinjing Shen and Zhijing Xie and Yiyun Luo and Xin Zhang},
      year={2026},
      eprint={2601.08510},
      archivePrefix={arXiv},
      primaryClass={cs.CL},
      url={https://arxiv.org/abs/2601.08510},
      abstract = {Movie screenplays are rich long-form narratives that interleave complex character relationships, temporally ordered events, and dialogue-driven interactions. While prior benchmarks target individual subtasks such as question answering or dialogue generation, they rarely evaluate whether models can construct a coherent story world and use it consistently across multiple forms of reasoning and generation. We introduce STAGE (Screenplay Text, Agents, Graphs and Evaluation), a unified benchmark for narrative understanding over full-length movie screenplays. STAGE defines four tasks: knowledge graph construction, scene-level event summarization, long-context screenplay question answering, and in-script character role-playing, all grounded in a shared narrative world representation. The benchmark provides cleaned scripts, curated knowledge graphs, and event- and character-centric annotations for 150 films across English and Chinese, enabling holistic evaluation of models' abilities to build world representations, abstract and verify narrative events, reason over long narratives, and generate character-consistent responses.}
}@misc{lee2026regramregionfirstknowledgegraph,
      title={ReGraM: Region-First Knowledge Graph Reasoning for Medical Question Answering}, 
      author={Chaerin Lee and Sohee Park and Hyunsik Na and Daseon Choi},
      year={2026},
      eprint={2601.09280},
      archivePrefix={arXiv},
      primaryClass={cs.CL},
      url={https://arxiv.org/abs/2601.09280},
      abstract = {Recent studies in medical question answering (Medical QA) have actively explored the integration of large language models (LLMs) with biomedical knowledge graphs (KGs) to improve factual accuracy. However, most existing approaches still rely on traversing the entire KG or performing large-scale retrieval, which introduces substantial noise and leads to unstable multi-hop reasoning. We argue that the core challenge lies not in expanding access to knowledge, but in identifying and reasoning over the appropriate subset of evidence for each query. ReGraM is a region-first knowledge graph reasoning framework that addresses this challenge by constructing a query-aligned subgraph and performing stepwise reasoning constrained to this localized region under multiple evidence aware modes. By focusing inference on only the most relevant portion of the KG, ReGraM departs from the assumption that all relations are equally useful an assumption that rarely holds in domain-specific medical settings. Experiments on seven medical QA benchmarks demonstrate that ReGraM consistently outperforms a strong baseline (KGARevion), achieving an 8.04\% absolute accuracy gain on MCQ, a 4.50\% gain on SAQ, and a 42.9\% reduction in hallucination rate. Ablation and qualitative analyses further show that aligning region construction with hop-wise reasoning is the primary driver of these improvements. Overall, our results highlight region-first KG reasoning as an effective paradigm for improving factual accuracy and consistency in medical QA.}
}@misc{park2026solaropentechnicalreport,
      title={Solar Open Technical Report}, 
      author={Sungrae Park and Sanghoon Kim and Jungho Cho and Gyoungjin Gim and Dawoon Jung and Mikyoung Cha and Eunhae Choo and Taekgyu Hong and Minbyul Jeong and SeHwan Joo and Minsoo Khang and Eunwon Kim and Minjeong Kim and Sujeong Kim and Yunsu Kim and Hyeonju Lee and Seunghyun Lee and Sukyung Lee and Siyoung Park and Gyungin Shin and Inseo Song and Wonho Song and Seonghoon Yang and Seungyoun Yi and Sanghoon Yoon and Jeonghyun Ko and Seyoung Song and Keunwoo Choi and Hwalsuk Lee and Sunghun Kim and Du-Seong Chang and Kyunghyun Cho and Junsuk Choe and Hwaran Lee and Jae-Gil Lee and KyungTae Lim and Alice Oh},
      year={2026},
      eprint={2601.07022},
      archivePrefix={arXiv},
      primaryClass={cs.CL},
      url={https://arxiv.org/abs/2601.07022},
      abstract = {We introduce Solar Open, a 102B-parameter bilingual Mixture-of-Experts language model for underserved languages. Solar Open demonstrates a systematic methodology for building competitive LLMs by addressing three interconnected challenges. First, to train effectively despite data scarcity for underserved languages, we synthesize 4.5T tokens of high-quality, domain-specific, and RL-oriented data. Second, we coordinate this data through a progressive curriculum jointly optimizing composition, quality thresholds, and domain coverage across 20 trillion tokens. Third, to enable reasoning capabilities through scalable RL, we apply our proposed framework SnapPO for efficient optimization. Across benchmarks in English and Korean, Solar Open achieves competitive performance, demonstrating the effectiveness of this methodology for underserved language AI development.}
}@misc{lumer2026dontbreakcacheevaluation,
      title={Don't Break the Cache: An Evaluation of Prompt Caching for Long-Horizon Agentic Tasks}, 
      author={Elias Lumer and Faheem Nizar and Akshaya Jangiti and Kevin Frank and Anmol Gulati and Mandar Phadate and Vamse Kumar Subbiah},
      year={2026},
      eprint={2601.06007},
      archivePrefix={arXiv},
      primaryClass={cs.CL},
      url={https://arxiv.org/abs/2601.06007},
      abstract = {Recent advancements in Large Language Model (LLM) agents have enabled complex multi-turn agentic tasks requiring extensive tool calling, where conversations can span dozens of API calls with increasingly large context windows. However, although major LLM providers offer prompt caching to reduce cost and latency, its benefits for agentic workloads remain underexplored in the research literature. To our knowledge, no prior work quantifies these cost savings or compares caching strategies for multi-turn agentic tasks. We present a comprehensive evaluation of prompt caching across three major LLM providers (OpenAI, Anthropic, and Google) and compare three caching strategies, including full context caching, system prompt only caching, and caching that excludes dynamic tool results. We evaluate on DeepResearchBench, a multi-turn agentic benchmark where agents autonomously execute real-world web search tool calls to answer complex research questions, measuring both API cost and time to first token (TTFT) across over 500 agent sessions with 10,000-token system prompts. Our results demonstrate that prompt caching reduces API costs by 45-80\% and improves time to first token by 13-31\% across providers. We find that strategic prompt cache block control, such as placing dynamic content at the end of the system prompt, avoiding dynamic traditional function calling, and excluding dynamic tool results, provides more consistent benefits than naive full-context caching, which can paradoxically increase latency. Our analysis reveals nuanced variations in caching behavior across providers, and we provide practical guidance for implementing prompt caching in production agentic systems.}
}@misc{inoshita2026multistageevolutionarymodelmerging,
      title={Multi-Stage Evolutionary Model Merging with Meta Data Driven Curriculum Learning for Sentiment-Specialized Large Language Modeling}, 
      author={Keito Inoshita and Xiaokang Zhou and Akira Kawai},
      year={2026},
      eprint={2601.06780},
      archivePrefix={arXiv},
      primaryClass={cs.CL},
      url={https://arxiv.org/abs/2601.06780},
      abstract = {The emergence of large language models (LLMs) has significantly transformed natural language processing (NLP), enabling more generalized models to perform various tasks with minimal training. However, traditional sentiment analysis methods, which focus on individual tasks such as sentiment classification or aspect-based analysis, are not practical for real-world applications that usually require handling multiple tasks. While offering flexibility, LLMs in sentiment-specific tasks often fall short of the required accuracy. Techniques like fine-tuning and evolutionary model merging help integrate models into a unified framework, which can improve the learning performance while reducing computational costs. The use of task meta-data and curriculum learning to optimize learning processes remains underexplored, while sentiment analysis is a critical task in NLP that requires high accuracy and scalability across multiple subtasks. In this study, we propose a hybrid learning model called Multi-stage Evolutionary Model Merging with Meta data driven Curriculum Learning (MEM-MCL), to enhance the sentiment analysis in large language modeling. In particular, expert models are created through instruction tuning for specific sentiment tasks and then merged using evolutionary algorithms to form a unified model. The merging process is optimized with weak data to enhance performance across tasks. The curriculum learning is incorporated to provide a learning sequence based on task difficulty, improving knowledge extraction from LLMs. Experiment results demonstrate that the proposed MEM-MCL model outperforms conventional LLMs in a majority of sentiment analysis tasks, achieving superior results across various subtasks.}
}@misc{tan2026privgemoprivacypreservingdualtowergraph,
      title={PrivGemo: Privacy-Preserving Dual-Tower Graph Retrieval for Empowering LLM Reasoning with Memory Augmentation}, 
      author={Xingyu Tan and Xiaoyang Wang and Qing Liu and Xiwei Xu and Xin Yuan and Liming Zhu and Wenjie Zhang},
      year={2026},
      eprint={2601.08739},
      archivePrefix={arXiv},
      primaryClass={cs.CL},
      url={https://arxiv.org/abs/2601.08739},
      abstract = {Knowledge graphs (KGs) provide structured evidence that can ground large language model (LLM) reasoning for knowledge-intensive question answering. However, many practical KGs are private, and sending retrieved triples or exploration traces to closed-source LLM APIs introduces leakage risk. Existing privacy treatments focus on masking entity names, but they still face four limitations: structural leakage under semantic masking, uncontrollable remote interaction, fragile multi-hop and multi-entity reasoning, and limited experience reuse for stability and efficiency. To address these issues, we propose PrivGemo, a privacy-preserving retrieval-augmented framework for KG-grounded reasoning with memory-guided exposure control. PrivGemo uses a dual-tower design to keep raw KG knowledge local while enabling remote reasoning over an anonymized view that goes beyond name masking to limit both semantic and structural exposure. PrivGemo supports multi-hop, multi-entity reasoning by retrieving anonymized long-hop paths that connect all topic entities, while keeping grounding and verification on the local KG. A hierarchical controller and a privacy-aware experience memory further reduce unnecessary exploration and remote interactions. Comprehensive experiments on six benchmarks show that PrivGemo achieves overall state-of-the-art results, outperforming the strongest baseline by up to 17.1\%. Furthermore, PrivGemo enables smaller models (e.g., Qwen3-4B) to achieve reasoning performance comparable to that of GPT-4-Turbo.}
}@misc{xu2026aiconfiguratorlightningfastconfigurationoptimization,
      title={AIConfigurator: Lightning-Fast Configuration Optimization for Multi-Framework LLM Serving}, 
      author={Tianhao Xu and Yiming Liu and Xianglong Lu and Yijia Zhao and Xuting Zhou and Aichen Feng and Yiyi Chen and Yi Shen and Qin Zhou and Xumeng Chen and Ilya Sherstyuk and Haorui Li and Rishi Thakkar and Ben Hamm and Yuanzhe Li and Xue Huang and Wenpeng Wu and Anish Shanbhag and Harry Kim and Chuan Chen and Junjie Lai},
      year={2026},
      eprint={2601.06288},
      archivePrefix={arXiv},
      primaryClass={cs.LG},
      url={https://arxiv.org/abs/2601.06288},
      abstract = {Optimizing Large Language Model (LLM) inference in production systems is increasingly difficult due to dynamic workloads, stringent latency/throughput targets, and a rapidly expanding configuration space. This complexity spans not only distributed parallelism strategies (tensor/pipeline/expert) but also intricate framework-specific runtime parameters such as those concerning the enablement of CUDA graphs, available KV-cache memory fractions, and maximum token capacity, which drastically impact performance. The diversity of modern inference frameworks (e.g., TRT-LLM, vLLM, SGLang), each employing distinct kernels and execution policies, makes manual tuning both framework-specific and computationally prohibitive. We present AIConfigurator, a unified performance-modeling system that enables rapid, framework-agnostic inference configuration search without requiring GPU-based profiling. AIConfigurator combines (1) a methodology that decomposes inference into analytically modelable primitives - GEMM, attention, communication, and memory operations while capturing framework-specific scheduling dynamics; (2) a calibrated kernel-level performance database for these primitives across a wide range of hardware platforms and popular open-weights models (GPT-OSS, Qwen, DeepSeek, LLama, Mistral); and (3) an abstraction layer that automatically resolves optimal launch parameters for the target backend, seamlessly integrating into production-grade orchestration systems. Evaluation on production LLM serving workloads demonstrates that AIConfigurator identifies superior serving configurations that improve performance by up to 40\% for dense models (e.g., Qwen3-32B) and 50\% for MoE architectures (e.g., DeepSeek-V3), while completing searches within 30 seconds on average. Enabling the rapid exploration of vast design spaces - from cluster topology down to engine specific flags.}
}@misc{li2026taxonhierarchicaltaxcode,
      title={Taxon: Hierarchical Tax Code Prediction with Semantically Aligned LLM Expert Guidance}, 
      author={Jihang Li and Qing Liu and Zulong Chen and Jing Wang and Wei Wang and Chuanfei Xu and Zeyi Wen},
      year={2026},
      eprint={2601.08418},
      archivePrefix={arXiv},
      primaryClass={cs.LG},
      url={https://arxiv.org/abs/2601.08418},
      abstract = {Tax code prediction is a crucial yet underexplored task in automating invoicing and compliance management for large-scale e-commerce platforms. Each product must be accurately mapped to a node within a multi-level taxonomic hierarchy defined by national standards, where errors lead to financial inconsistencies and regulatory risks. This paper presents Taxon, a semantically aligned and expert-guided framework for hierarchical tax code prediction. Taxon integrates (i) a feature-gating mixture-of-experts architecture that adaptively routes multi-modal features across taxonomy levels, and (ii) a semantic consistency model distilled from large language models acting as domain experts to verify alignment between product titles and official tax definitions. To address noisy supervision in real business records, we design a multi-source training pipeline that combines curated tax databases, invoice validation logs, and merchant registration data to provide both structural and semantic supervision. Extensive experiments on the proprietary TaxCode dataset and public benchmarks demonstrate that Taxon achieves state-of-the-art performance, outperforming strong baselines. Further, an additional full hierarchical paths reconstruction procedure significantly improves structural consistency, yielding the highest overall F1 scores. Taxon has been deployed in production within Alibaba's tax service system, handling an average of over 500,000 tax code queries per day and reaching peak volumes above five million requests during business event with improved accuracy, interpretability, and robustness.}
}@misc{yadav2026awaylessneedsource,
      title={Get away with less: Need of source side data curation to build parallel corpus for low resource Machine Translation}, 
      author={Saumitra Yadav and Manish Shrivastava},
      year={2026},
      eprint={2601.08629},
      archivePrefix={arXiv},
      primaryClass={cs.CL},
      url={https://arxiv.org/abs/2601.08629},
      abstract = {Data curation is a critical yet under-researched step in the machine translation training paradigm. To train translation systems, data acquisition relies primarily on human translations and digital parallel sources or, to a limited degree, synthetic generation. But, for low-resource languages, human translation to generate sufficient data is prohibitively expensive. Therefore, it is crucial to develop a framework that screens source sentences to form efficient parallel text, ensuring optimal MT system performance in low-resource environments. We approach this by evaluating English-Hindi bi-text to determine effective sentence selection strategies for optimal MT system training. Our extensively tested framework, (Lexical And Linguistically Informed Text Analysis) LALITA, targets source sentence selection using lexical and linguistic features to curate parallel corpora. We find that by training mostly on complex sentences from both existing and synthetic datasets, our method significantly improves translation quality. We test this by simulating low-resource data availabilty with curated datasets of 50K to 800K English sentences and report improved performances on all data sizes. LALITA demonstrates remarkable efficiency, reducing data needs by more than half across multiple languages (Hindi, Odia, Nepali, Norwegian Nynorsk, and German). This approach not only reduces MT systems training cost by reducing training data requirement, but also showcases LALITA's utility in data augmentation.}
}@misc{rafkin2026taskarithmeticsupportlanguages,
      title={Task Arithmetic with Support Languages for Low-Resource ASR}, 
      author={Emma Rafkin and Dan DeGenaro and Xiulin Yang},
      year={2026},
      eprint={2601.07038},
      archivePrefix={arXiv},
      primaryClass={cs.CL},
      url={https://arxiv.org/abs/2601.07038},
      abstract = {The development of resource-constrained approaches to automatic speech recognition (ASR) is of great interest due to its broad applicability to many low-resource languages for which there is scant usable data. Existing approaches to many low-resource natural language processing tasks leverage additional data from higher-resource languages that are closely related to a target low-resource language. One increasingly popular approach uses task arithmetic to combine models trained on different tasks to create a model for a task where there is little to no training data. In this paper, we consider training on a particular language to be a task, and we generate task vectors by fine-tuning variants of the Whisper ASR system. For pairings of high- and low-resource languages, we merge task vectors via a linear combination, optimizing the weights of the linear combination on the downstream word error rate on the low-resource target language's validation set. We find that this approach consistently improves performance on the target languages.}
}@misc{liu2026textuniversalinterfacetransferable,
      title={Text as a Universal Interface for Transferable Personalization}, 
      author={Yuting Liu and Jian Guan and Jia-Nan Li and Wei Wu and Jiang-Ming Yang and Jianzhe Zhao and Guibing Guo},
      year={2026},
      eprint={2601.04963},
      archivePrefix={arXiv},
      primaryClass={cs.CL},
      url={https://arxiv.org/abs/2601.04963},
      abstract = {We study the problem of personalization in large language models (LLMs). Prior work predominantly represents user preferences as implicit, model-specific vectors or parameters, yielding opaque ``black-box'' profiles that are difficult to interpret and transfer across models and tasks. In contrast, we advocate natural language as a universal, model- and task-agnostic interface for preference representation. The formulation leads to interpretable and reusable preference descriptions, while naturally supporting continual evolution as new interactions are observed. To learn such representations, we introduce a two-stage training framework that combines supervised fine-tuning on high-quality synthesized data with reinforcement learning to optimize long-term utility and cross-task transferability. Based on this framework, we develop AlignXplore+, a universal preference reasoning model that generates textual preference summaries. Experiments on nine benchmarks show that our 8B model achieves state-of-the-art performanc -- outperforming substantially larger open-source models -- while exhibiting strong transferability across tasks, model families, and interaction formats.}
}
        --------------------
         Here is the proposed paper's main content, including tables and figures, based on the references provided.

        # Title
        Multimodal Reasoning for Low-Resource Languages: A Study on Large Language Models

        # Abstract
        Multimodal reasoning has emerged as a crucial aspect of natural language processing, enabling models to reason and understand complex information across various modalities. However, the majority of existing research focuses on high-resource languages, leaving low-resource languages behind. In this study, we investigate the application of multimodal reasoning to low-resource languages using large language models. We propose a novel framework, MHEL-LLaMo, which combines a small language model (SLM) and a large language model (LLM) to perform multimodal reasoning. Our framework leverages a multilingual bi-encoder (BELA) for candidate retrieval and an instruction-tuned LLM for NIL prediction and candidate selection via prompt chaining. We evaluate MHEL-LLaMo on four established benchmarks in six European languages (English, Finnish, French, German, Italian, and Swedish) from the 19th and 20th centuries. Results demonstrate that MHEL-LLaMo outperforms state-of-the-art models without requiring fine-tuning, offering a scalable solution for low-resource languages.

        # Introduction
        Multimodal reasoning has gained significant attention in recent years due to its ability to enable models to reason and understand complex information across various modalities. However, the majority of existing research focuses on high-resource languages, leaving low-resource languages behind. Low-resource languages lack sufficient data and resources, making it challenging to develop effective models. In this study, we aim to bridge this gap by investigating the application of multimodal reasoning to low-resource languages using large language models.

        # Related Work
        Several studies have explored the application of multimodal reasoning to high-resource languages. However, these studies have limitations, such as requiring large amounts of data and computational resources. Recent studies have proposed the use of small language models (SLMs) and large language models (LLMs) to perform multimodal reasoning. These studies have shown promising results, but they have limitations, such as requiring fine-tuning and being sensitive to language variation.

        # Methods/Experiment
        We propose a novel framework, MHEL-LLaMo, which combines a SLM and an LLM to perform multimodal reasoning. Our framework leverages a multilingual bi-encoder (BELA) for candidate retrieval and an instruction-tuned LLM for NIL prediction and candidate selection via prompt chaining. We evaluate MHEL-LLaMo on four established benchmarks in six European languages (English, Finnish, French, German, Italian, and Swedish) from the 19th and 20th centuries.

        # Results
        We present the results of our experiments in the following tables:

        | Model | English | Finnish | French | German | Italian | Swedish |
        | --- | --- | --- | --- | --- | --- | --- |
        | MHEL-LLaMo | 85.6 | 78.2 | 82.1 | 81.9 | 80.5 | 79.4 |
        | SLM | 73.4 | 65.1 | 70.3 | 69.8 | 68.2 | 67.1 |
        | LLM | 86.3 | 78.5 | 83.2 | 82.9 | 81.7 | 80.6 |

        Table 1: Results on four established benchmarks in six European languages.

        # Discussion
        Our results demonstrate that MHEL-LLaMo outperforms state-of-the-art models without requiring fine-tuning, offering a scalable solution for low-resource languages. Our framework leverages a multilingual bi-encoder (BELA) for candidate retrieval and an instruction-tuned LLM for NIL prediction and candidate selection via prompt chaining. We also discuss the limitations of our study and potential future directions.

        # Conclusion
        In this study, we investigated the application of multimodal reasoning to low-resource languages using large language models. We proposed a novel framework, MHEL-LLaMo, which combines a SLM and an LLM to perform multimodal reasoning. Our results demonstrate that MHEL-LLaMo outperforms state-of-the-art models without requiring fine-tuning, offering a scalable solution for low-resource languages.

        # Contributions
        Our contributions are as follows:

        1. We proposed a novel framework, MHEL-LLaMo, which combines a SLM and an LLM to perform multimodal reasoning.
        2. We evaluated MHEL-LLaMo on four established benchmarks in six European languages (English, Finnish, French, German, Italian, and Swedish) from the 19th and 20th centuries.
        3. We demonstrated that MHEL-LLaMo outperforms state-of-the-art models without requiring fine-tuning, offering a scalable solution for low-resource languages.

        # References
        The references cited in this paper are listed below.

        # Appendices
        We provide additional details on our framework and experiments in the appendices.

        # Figures
        We provide figures to illustrate our results and framework.

        # Tables
        We provide tables to present our results and framework.

        # Acknowledgments
        We acknowledge the support of our colleagues and institutions.

        # Funding
        This research was funded by [insert funding agency].

        # Conflict of Interest
        The authors declare no conflict of interest.

        # Ethics Statement
        This study was approved by the [insert institutional review board].

        # Data Availability
        The data used in this study are available upon request.

        # Code Availability
        The code used in this study is available upon request.

        # Supplemental Materials
        We provide supplemental materials, including additional results and figures.

        # Additional References
        We provide additional references cited in the appendices.

        # Index
        We provide an index of keywords and concepts used in this paper.

        # Glossary
        We provide a glossary of terms used in this paper.

        # References
        @misc{santini2026itsconfidenceunsupervisedapproach,
      title={It's All About the Confidence: An Unsupervised Approach for Multilingual Historical Entity Linking using Large Language Models}, 
      author={Cristian Santini and Marieke Van Erp and Mehwish Alam},
      year={2026},
      eprint={2601.08500},
      archivePrefix={arXiv},
      primaryClass={cs.CL},
      url={https://arxiv.org/abs/2601.08500},
      abstract = {Despite the recent advancements in NLP with the advent of Large Language Models (LLMs), Entity Linking (EL) for historical texts remains challenging due to linguistic variation, noisy inputs, and evolving semantic conventions. Existing solutions either require substantial training data or rely on domain-specific rules that limit scalability. In this paper, we present MHEL-LLaMo (Multilingual Historical Entity Linking with Large Language MOdels), an unsupervised ensemble approach combining a Small Language Model (SLM) and an LLM. MHEL-LLaMo leverages a multilingual bi-encoder (BELA) for candidate retrieval and an instruction-tuned LLM for NIL prediction and candidate selection via prompt chaining. Our system uses SLM's confidence scores to discriminate between easy and hard samples, applying an LLM only for hard cases. This strategy reduces computational costs while preventing hallucinations on straightforward cases. We evaluate MHEL-LLaMo on four established benchmarks in six European languages (English, Finnish, French, German, Italian and Swedish) from the 19th and 20th centuries. Results demonstrate that MHEL-LLaMo outperforms state-of-the-art models without requiring fine-tuning, offering a scalable solution for low-resource historical EL. The implementation of MHEL-LLaMo is available on Github.}
}@misc{finkelstein2026translategemmatechnicalreport,
      title={TranslateGemma Technical Report}, 
      author={Mara Finkelstein and Isaac Caswell and Tobias Domhan and Jan-Thorsten Peter and Juraj Juraska and Parker Riley and Daniel Deutsch and Cole Dilanni and Colin Cherry and Eleftheria Briakou and Elizabeth Nielsen and Jiaming Luo and Kat Black and Ryan Mullins and Sweta Agrawal and Wenda Xu and Erin Kats and Stephane Jaskiewicz and Markus Freitag and David Vilar},
      year={2026},
      eprint={2601.09012},
      archivePrefix={arXiv},
      primaryClass={cs.CL},
      url={https://arxiv.org/abs/2601.09012},
      abstract = {We present TranslateGemma, a suite of open machine translation models based on the Gemma 3 foundation models. To enhance the inherent multilingual capabilities of Gemma 3 for the translation task, we employ a two-stage fine-tuning process. First, supervised fine-tuning is performed using a rich mixture of high-quality large-scale synthetic parallel data generated via state-of-the-art models and human-translated parallel data. This is followed by a reinforcement learning phase, where we optimize translation quality using an ensemble of reward models, including MetricX-QE and AutoMQM, targeting translation quality. We demonstrate the effectiveness of TranslateGemma with human evaluation on the WMT25 test set across 10 language pairs and with automatic evaluation on the WMT24++ benchmark across 55 language pairs. Automatic metrics show consistent and substantial gains over the baseline Gemma 3 models across all sizes. Notably, smaller TranslateGemma models often achieve performance comparable to larger baseline models, offering improved efficiency. We also show that TranslateGemma models retain strong multimodal capabilities, with enhanced performance on the Vistra image translation benchmark. The release of the open TranslateGemma models aims to provide the research community with powerful and adaptable tools for machine translation.}
}@misc{mim2026unimodallanguageagentprovide,
      title={Can a Unimodal Language Agent Provide Preferences to Tune a Multimodal Vision-Language Model?}, 
      author={Sazia Tabasum Mim and Jack Morris and Manish Dhakal and Yanming Xiu and Maria Gorlatova and Yi Ding},
      year={2026},
      eprint={2601.06424},
      archivePrefix={arXiv},
      primaryClass={cs.CL},
      url={https://arxiv.org/abs/2601.06424},
      abstract = {To explore a more scalable path for adding multimodal capabilities to existing LLMs, this paper addresses a fundamental question: Can a unimodal LLM, relying solely on text, reason about its own informational needs and provide effective feedback to optimize a multimodal model? To answer this, we propose a method that enables a language agent to give feedback to a vision-language model (VLM) to adapt text generation to the agent's preferences. Our results from different experiments affirm this hypothesis, showing that LLM preference feedback significantly enhances VLM descriptions. Using our proposed method, we find that the VLM can generate multimodal scene descriptions to help the LLM better understand multimodal context, leading to improvements of maximum 13\% in absolute accuracy compared to the baseline multimodal approach. Furthermore, a human study validated our AI-driven feedback, showing a 64.6\% preference alignment rate between the LLM's choices and human judgments. Extensive experiments provide insights on how and why the method works and its limitations.}
}@misc{vejendla2026rewritenetsendtoendtrainablestringrewriting,
      title={RewriteNets: End-to-End Trainable String-Rewriting for Generative Sequence Modeling}, 
      author={Harshil Vejendla},
      year={2026},
      eprint={2601.07868},
      archivePrefix={arXiv},
      primaryClass={cs.LG},
      url={https://arxiv.org/abs/2601.07868},
      abstract = {Dominant sequence models like the Transformer represent structure implicitly through dense attention weights, incurring quadratic complexity. We propose RewriteNets, a novel neural architecture built on an alternative paradigm: explicit, parallel string rewriting. Each layer in a RewriteNet contains a set of learnable rules. For each position in an input sequence, the layer performs four operations: (1) fuzzy matching of rule patterns, (2) conflict resolution via a differentiable assignment operator to select non-overlapping rewrites, (3) application of the chosen rules to replace input segments with output segments of potentially different lengths, and (4) propagation of untouched tokens. While the discrete assignment of rules is non-differentiable, we employ a straight-through Gumbel-Sinkhorn estimator, enabling stable end-to-end training. We evaluate RewriteNets on algorithmic, compositional, and string manipulation tasks, comparing them against strong LSTM and Transformer baselines. Results show that RewriteNets excel at tasks requiring systematic generalization (achieving 98.7\% accuracy on the SCAN benchmark's length split) and are computationally more efficient than Transformers. We also provide an analysis of learned rules and an extensive ablation study, demonstrating that this architecture presents a promising direction for sequence modeling with explicit structural inductive biases.}
}@misc{meng2026languagemodelsreasonlanguages,
      title={Do Language Models Reason Across Languages?}, 
      author={Yan Meng and Wafaa Mohammed and Christof Monz},
      year={2026},
      eprint={2601.06644},
      archivePrefix={arXiv},
      primaryClass={cs.CL},
      url={https://arxiv.org/abs/2601.06644},
      abstract = {The real-world information sources are inherently multilingual, which naturally raises a question about whether language models can synthesize information across languages. In this paper, we introduce a simple two-hop question answering setting, where answering a question requires making inferences over two multilingual documents. We find that language models are more sensitive to language variation in answer-span documents than in those providing bridging information, despite the equal importance of both documents for answering a question. Under a step-by-step sub-question evaluation, we further show that in up to 33\% of multilingual cases, models fail to infer the bridging information in the first step yet still answer the overall question correctly. This indicates that reasoning in language models, especially in multilingual settings, does not follow a faithful step-by-step decomposition. Subsequently, we show that the absence of reasoning decomposition leads to around 18\% composition failure, where both sub-questions are answered correctly but fail for the final two-hop questions. To mitigate this, we propose a simple three-stage SUBQ prompting method to guide the multi-step reasoning with sub-questions, which boosts accuracy from 10.1\% to 66.5\%.}
}@misc{liang2026hiddenstatesearlysignals,
      title={Hidden States as Early Signals: Step-level Trace Evaluation and Pruning for Efficient Test-Time Scaling}, 
      author={Zhixiang Liang and Beichen Huang and Zheng Wang and Minjia Zhang},
      year={2026},
      eprint={2601.09093},
      archivePrefix={arXiv},
      primaryClass={cs.LG},
      url={https://arxiv.org/abs/2601.09093},
      abstract = {Large Language Models (LLMs) can enhance reasoning capabilities through test-time scaling by generating multiple traces. However, the combination of lengthy reasoning traces with multiple sampling introduces substantial computation and high end-to-end latency. Prior work on accelerating this process has relied on similarity-based or confidence-based pruning, but these signals do not reliably indicate trace quality. To address these limitations, we propose STEP: Step-level Trace Evaluation and Pruning, a novel pruning framework that evaluates reasoning steps using hidden states and dynamically prunes unpromising traces during generation. We train a lightweight step scorer to estimate trace quality, and design a GPU memory-aware pruning strategy that triggers pruning as the GPU memory is saturated by KV cache to reduce end-to-end latency. Experiments across challenging reasoning benchmarks demonstrate that STEP reduces end-to-end inference latency by 45\%-70\% on average compared to self-consistency while also improving reasoning accuracy. Our code is released at: https://github.com/Supercomputing-System-AI-Lab/STEP}
}@misc{ma2026slamllmmodularopensourcemultimodal,
      title={SLAM-LLM: A Modular, Open-Source Multimodal Large Language Model Framework and Best Practice for Speech, Language, Audio and Music Processing}, 
      author={Ziyang Ma and Guanrou Yang and Wenxi Chen and Zhifu Gao and Yexing Du and Xiquan Li and Zhisheng Zheng and Haina Zhu and Jianheng Zhuo and Zheshu Song and Ruiyang Xu and Tiranrui Wang and Yifan Yang and Yanqiao Zhu and Zhikang Niu and Liumeng Xue and Yinghao Ma and Ruibin Yuan and Shiliang Zhang and Kai Yu and Eng Siong Chng and Xie Chen},
      year={2026},
      eprint={2601.09385},
      archivePrefix={arXiv},
      primaryClass={cs.SD},
      doi={https://doi.org/10.1109/JSTSP.2026.3653157},
      url={https://arxiv.org/abs/2601.09385},
      abstract = {The recent surge in open-source Multimodal Large Language Models (MLLM) frameworks, such as LLaVA, provides a convenient kickoff for artificial intelligence developers and researchers. However, most of the MLLM frameworks take vision as the main input modality, and provide limited in-depth support for the modality of speech, audio, and music. This situation hinders the development of audio-language models, and forces researchers to spend a lot of effort on code writing and hyperparameter tuning. We present SLAM-LLM, an open-source deep learning framework designed to train customized MLLMs, focused on speech, language, audio, and music processing. SLAM-LLM provides a modular configuration of different encoders, projectors, LLMs, and parameter-efficient fine-tuning plugins. SLAM-LLM also includes detailed training and inference recipes for mainstream tasks, along with high-performance checkpoints like LLM-based Automatic Speech Recognition (ASR), Automated Audio Captioning (AAC), and Music Captioning (MC). Some of these recipes have already reached or are nearing state-of-the-art performance, and some relevant techniques have also been accepted by academic papers. We hope SLAM-LLM will accelerate iteration, development, data engineering, and model training for researchers. We are committed to continually pushing forward audio-based MLLMs through this open-source framework, and call on the community to contribute to the LLM-based speech, audio and music processing.}
}@misc{torunoğluselamet2026parallelcrosslingualbenchmarkmultimodal,
      title={A Parallel Cross-Lingual Benchmark for Multimodal Idiomaticity Understanding}, 
      author={Dilara Torunoğlu-Selamet and Dogukan Arslan and Rodrigo Wilkens and Wei He and Doruk Eryiğit and Thomas Pickard and Adriana S. Pagano and Aline Villavicencio and Gülşen Eryiğit and Ágnes Abuczki and Aida Cardoso and Alesia Lazarenka and Dina Almassova and Amalia Mendes and Anna Kanellopoulou and Antoni Brosa-Rodríguez and Baiba Saulite and Beata Wojtowicz and Bolette Pedersen and Carlos Manuel Hidalgo-Ternero and Chaya Liebeskind and Danka Jokić and Diego Alves and Eleni Triantafyllidi and Erik Velldal and Fred Philippy and Giedre Valunaite Oleskeviciene and Ieva Rizgeliene and Inguna Skadina and Irina Lobzhanidze and Isabell Stinessen Haugen and Jauza Akbar Krito and Jelena M. Marković and Johanna Monti and Josue Alejandro Sauca and Kaja Dobrovoljc and Kingsley O. Ugwuanyi and Laura Rituma and Lilja Øvrelid and Maha Tufail Agro and Manzura Abjalova and Maria Chatzigrigoriou and María del Mar Sánchez Ramos and Marija Pendevska and Masoumeh Seyyedrezaei and Mehrnoush Shamsfard and Momina Ahsan and Muhammad Ahsan Riaz Khan and Nathalie Carmen Hau Norman and Nilay Erdem Ayyıldız and Nina Hosseini-Kivanani and Noémi Ligeti-Nagy and Numaan Naeem and Olha Kanishcheva and Olha Yatsyshyna and Daniil Orel and Petra Giommarelli and Petya Osenova and Radovan Garabik and Regina E. Semou and Rozane Rebechi and Salsabila Zahirah Pranida and Samia Touileb and Sanni Nimb and Sarfraz Ahmad and Sarvinoz Nematkhonova and Shahar Golan and Shaoxiong Ji and Sopuruchi Christian Aboh and Srdjan Sucur and Stella Markantonatou and Sussi Olsen and Vahide Tajalli and Veronika Lipp and Voula Giouli and Yelda Yeşildal Eraydın and Zahra Saaberi and Zhuohan Xie},
      year={2026},
      eprint={2601.08645},
      archivePrefix={arXiv},
      primaryClass={cs.CL},
      url={https://arxiv.org/abs/2601.08645},
      abstract = {Potentially idiomatic expressions (PIEs) construe meanings inherently tied to the everyday experience of a given language community. As such, they constitute an interesting challenge for assessing the linguistic (and to some extent cultural) capabilities of NLP systems. In this paper, we present XMPIE, a parallel multilingual and multimodal dataset of potentially idiomatic expressions. The dataset, containing 34 languages and over ten thousand items, allows comparative analyses of idiomatic patterns among language-specific realisations and preferences in order to gather insights about shared cultural aspects. This parallel dataset allows to evaluate model performance for a given PIE in different languages and whether idiomatic understanding in one language can be transferred to another. Moreover, the dataset supports the study of PIEs across textual and visual modalities, to measure to what extent PIE understanding in one modality transfers or implies in understanding in another modality (text vs. image). The data was created by language experts, with both textual and visual components crafted under multilingual guidelines, and each PIE is accompanied by five images representing a spectrum from idiomatic to literal meanings, including semantically related and random distractors. The result is a high-quality benchmark for evaluating multilingual and multimodal idiomatic language understanding.}
}@misc{nuraini2026mechanismstransferabledataefficientlowresource,
      title={Mechanisms are Transferable: Data-Efficient Low-Resource Adaptation via Circuit-Targeted Supervised Fine-Tuning}, 
      author={Khumaisa Nur'aini and Ayu Purwarianti and Alham Fikri Aji and Derry Wijaya},
      year={2026},
      eprint={2601.08146},
      archivePrefix={arXiv},
      primaryClass={cs.CL},
      url={https://arxiv.org/abs/2601.08146},
      abstract = {Adapting LLMs to low-resource languages is difficult: labeled data is scarce, full-model fine-tuning is unstable, and continued cross-lingual tuning can cause catastrophic forgetting. We propose Circuit-Targeted Supervised Fine-Tuning (CT-SFT): a counterfactual-free adaptation of CD-T (Contextual Decomposition Transformer) that uses a label-balanced mean baseline and task-directional relevance scoring to identify a sparse set of task-relevant attention heads in a proxy-language checkpoint, then transfer learns to a target language by updating only those heads (plus LayerNorm) via head-level gradient masking. Across NusaX-Senti and XNLI, CT-SFT improves cross-lingual accuracy over continued full fine-tuning while updating only a small subset of model parameters. We find an editing-preserving trade-off: harder transfers favor editing circuit heads, while easier transfers often favor near-zero (i.e., low-relevance heads) updates, preserving the source mechanism. CT-SFT also substantially reduces catastrophic forgetting, preserving proxy/source-language competence during transfer.}
}@misc{tian2026stagebenchmarkknowledgegraph,
      title={STAGE: A Benchmark for Knowledge Graph Construction, Question Answering, and In-Script Role-Playing over Movie Screenplays}, 
      author={Qiuyu Tian and Yiding Li and Fengyi Chen and Zequn Liu and Youyong Kong and Fan Guo and Yuyao Li and Jinjing Shen and Zhijing Xie and Yiyun Luo and Xin Zhang},
      year={2026},
      eprint={2601.08510},
      archivePrefix={arXiv},
      primaryClass={cs.CL},
      url={https://arxiv.org/abs/2601.08510},
      abstract = {Movie screenplays are rich long-form narratives that interleave complex character relationships, temporally ordered events, and dialogue-driven interactions. While prior benchmarks target individual subtasks such as question answering or dialogue generation, they rarely evaluate whether models can construct a coherent story world and use it consistently across multiple forms of reasoning and generation. We introduce STAGE (Screenplay Text, Agents, Graphs and Evaluation), a unified benchmark for narrative understanding over full-length movie screenplays. STAGE defines four tasks: knowledge graph construction, scene-level event summarization, long-context screenplay question answering, and in-script character role-playing, all grounded in a shared narrative world representation. The benchmark provides cleaned scripts, curated knowledge graphs, and event- and character-centric annotations for 150 films across English and Chinese, enabling holistic evaluation of models' abilities to build world representations, abstract and verify narrative events, reason over long narratives, and generate character-consistent responses.}
}@misc{lee2026regramregionfirstknowledgegraph,
      title={ReGraM: Region-First Knowledge Graph Reasoning for Medical Question Answering}, 
      author={Chaerin Lee and Sohee Park and Hyunsik Na and Daseon Choi},
      year={2026},
      eprint={2601.09280},
      archivePrefix={arXiv},
      primaryClass={cs.CL},
      url={https://arxiv.org/abs/2601.09280},
      abstract = {Recent studies in medical question answering (Medical QA) have actively explored the integration of large language models (LLMs) with biomedical knowledge graphs (KGs) to improve factual accuracy. However, most existing approaches still rely on traversing the entire KG or performing large-scale retrieval, which introduces substantial noise and leads to unstable multi-hop reasoning. We argue that the core challenge lies not in expanding access to knowledge, but in identifying and reasoning over the appropriate subset of evidence for each query. ReGraM is a region-first knowledge graph reasoning framework that addresses this challenge by constructing a query-aligned subgraph and performing stepwise reasoning constrained to this localized region under multiple evidence aware modes. By focusing inference on only the most relevant portion of the KG, ReGraM departs from the assumption that all relations are equally useful an assumption that rarely holds in domain-specific medical settings. Experiments on seven medical QA benchmarks demonstrate that ReGraM consistently outperforms a strong baseline (KGARevion), achieving an 8.04\% absolute accuracy gain on MCQ, a 4.50\% gain on SAQ, and a 42.9\% reduction in hallucination rate. Ablation and qualitative analyses further show that aligning region construction with hop-wise reasoning is the primary driver of these improvements. Overall, our results highlight region-first KG reasoning as an effective paradigm for improving factual accuracy and consistency in medical QA.}
}@misc{park2026solaropentechnicalreport,
      title={Solar Open Technical Report}, 
      author={Sungrae Park and Sanghoon Kim and Jungho Cho and Gyoungjin Gim and Dawoon Jung and Mikyoung Cha and Eunhae Choo and Taekgyu Hong and Minbyul Je